\documentclass[11pt]{article}
\usepackage[a4paper,margin=25mm]{geometry}
\usepackage{amsmath,amssymb,amsthm}
\usepackage[T1]{fontenc}
\usepackage{hyperref}
\hypersetup{colorlinks=true,linkcolor=blue,urlcolor=blue}
\newtheorem{theorem}{Theorem}
\newtheorem{corollary}{Corollary}
\newcommand{\KB}{\mathit{KB}}
\newcommand{\Dall}{D_{\mathrm{all}}}
\newcommand{\naf}{\mathop{\mathrm{not}}}
\newcommand{\den}[1]{\lbrack\!\lbrack #1\rbrack\!\rbrack_{M,\rho}}
\title{Soundness of an Untyped Anthill KB Logic}
\author{}
\date{}
\begin{document}
\maketitle

\section{Objective and result}
The objective is to prove that the logic preserves truth: a formula derived from
a knowledge base and local assumptions is true in any admissible model where
those assumptions hold. The result concerns derivations, independently of how
a derivation is discovered. No automatic search algorithm or native matcher is
needed for the logic-soundness theorem.

We define the syntax and inference rules, give their compositional semantics,
and prove soundness by induction on derivations. Induction principles are
ordinary KB clauses whose validity is established in the chosen model.
Negation as failure (NAF) requires independently justified evidence; the
formalized theorem does not identify arbitrary failed search with negation.
Existence of a model satisfying these conditions then establishes consistency.

The core theorem is \texttt{Semantics.Model.sound}. For the certified NAF
provider and inductively interpreted KB, its packaged form is:
\begin{quote}
\href{AnthillModKbSpec/Resolution/InductiveSoundness.lean}{\texttt{InductiveSoundness.declarative\_sound}}
\end{quote}
Proof-discovery experiments are described in the separate
\href{docs/implementation.md}{implementation note}.

\section{Formula syntax and KB clauses}
The shared syntax is \texttt{Term}, defined in \texttt{Syntax.lean}.
Writing $G$ for a goal or term and $H$ for a symbol-headed application, its grammar is:
\[
\begin{aligned}
 H ::= {}& f(G_1,\ldots,G_n)\qquad(n\geq 0),\\
 G ::= {}& v_i \mid c \mid H
       \mid \operatorname{dictionaryCall}(G,\pi,f,\bar G)\\
       &\mid \operatorname{equal}(G,G) \mid \top \mid \bot\\
       &\mid G\land G \mid G\lor G \mid G\Rightarrow G
        \mid \naf G \mid \forall v_i.G \mid \exists v_i.G.
\end{aligned}
\]
Here $v_i$ ranges over variable names, $c$ over literals, $f$ over symbols, and
$\pi$ over lists of dictionary field selectors. $H$ denotes an application;
its use as a predicate or an operation depends on its position and interpretation.
Data and formulas share this carrier, so application arguments are also $G$s.
Logical variables and predicate variables use the same variable constructor.

A predicate symbol as a value is a nullary application. The designated ordinary
symbol \texttt{ho\_apply} represents explicit predicate application. For a
concrete symbol value $p()$, the calculus relates
$\operatorname{ho\_apply}(p(),\bar t)$ to the ordinary atom $p(\bar t)$.
This does not introduce a separate kind of induction goal.

A KB is a set of clauses. The following forms use this same set:
\begin{center}
\begin{tabular}{ll}
\textbf{Form} & \textbf{Clause notation} \\
Predicate rule & $p(\bar t)\mathbin{:-}G_1,\ldots,G_k$ \\
Equation rule & $l\mathrel{\Leftrightarrow}r\mathbin{:-}G_1,\ldots,G_k$ \\
Goal clause (denial) & $\mathbin{:-}G_1,\ldots,G_k$
\end{tabular}
\end{center}
Formally, \texttt{HornClause} has two constructors:
\[
 \operatorname{definite}(G_0,[G_1,\ldots,G_k])
 \qquad\text{or}\qquad
 \operatorname{goal}([G_1,\ldots,G_k]).
\]
The first denotes $(G_1\land\cdots\land G_k)\Rightarrow G_0$; the second has
$\bot$ as its conclusion. An empty body denotes truth: a definite clause then
gives a fact or an unconditional equation, while an empty goal clause asserts
falsity. Free clause variables are implicitly universally quantified together.
Local assumptions $\Gamma$ instead have fixed free parameters.

Each body goal $G_i$ can be a predicate application, an equation, truth or falsity,
a conjunction or disjunction, a nested implication, NAF, or a quantified formula.
Operation expressions in these goals can contain dictionary calls and further
operation arguments. For example, both of these are clause forms:
\[
\begin{aligned}
 p(x)&\mathbin{:-}\operatorname{equal}(f(x),g(x)),\ q(x),\\
 h(x)\mathrel{\Leftrightarrow}k(x)
      &\mathbin{:-}\operatorname{isNat}(x),\ \operatorname{equal}(f(x),g(x)).
\end{aligned}
\]
Thus equations occur in heads and in bodies; they are goals to prove wherever
they occur in a body. Nested implication introduces local assumptions, quantifiers
bind their bodies, and NAF requires its separate justified evidence.

The Lean syntax is raw: \texttt{Atom} and \texttt{Formula} are aliases of
\texttt{Term}, so a definite head $G_0$ and each body goal are arbitrary terms at
the syntax level. Predicate and equality heads are common forms, rather than
an enforced head restriction. The model must validate every included clause;
the individual search fragments impose their own additional grammar restrictions.

A conditional operation equation $l\mathrel{\Leftrightarrow}r\mathbin{:-}\bar C$
is a definite clause with head $\operatorname{equal}(l,r)$, represented by
\texttt{HornClause.equation}. Thus predicate and equation rules occupy the same
KB set. It asserts equality of returned values whenever its condition body holds:
\[
 \forall\rho.\;M,\rho\models\bigwedge\bar C\quad\Longrightarrow\quad
 \lbrack\!\lbrack l\rbrack\!\rbrack_{M,\rho}
 =\lbrack\!\lbrack r\rbrack\!\rbrack_{M,\rho}.
\]
Both sides can contain dictionary calls. The same clause substitution
instantiates both sides and the body; using the equation requires proving that body.
Equality goals can also occur in the body of a predicate or equation rule;
proving them from other equation rules can therefore enable a rule's conclusion.

\section{The derivation rules}
The judgment
\[
 \KB;\Gamma\vdash_N G
\]
is implemented by \texttt{Derives} in \texttt{Derivation.lean}. The parameter
$N$ specifies which NAF evidence is available. With NAF disabled, it is the
positive calculus. A derivation is a finite tree built from the inference rules;
its definition does not specify clause-selection order or backtracking.

The calculus has the usual assumption, truth, conjunction, disjunction,
implication, falsity-elimination, and quantifier rules. In particular,
\[
 \frac{\KB;\Gamma,A\vdash_N B}{\KB;\Gamma\vdash_N A\Rightarrow B},
 \qquad
 \frac{\KB;\Gamma\vdash_N F[a/x]}{\KB;\Gamma\vdash_N\forall x.F(x)}.
\]
In the second rule, $a$ is a fresh eigenparameter, absent from the free variables
of $\Gamma$ and of the quantified goal. Opening is capture-avoiding.

For clause application, choose one capture-avoiding substitution $\sigma$ for
both head and body:
\[
 \frac{(G_0\mathbin{:-}\bar G)\in\KB\qquad
       \KB;\Gamma\vdash_N\sigma(\bigwedge\bar G)}
      {\KB;\Gamma\vdash_N\sigma(G_0)}.
\]
The rule does not substitute through $\Gamma$ or any surrounding goal. A goal
clause similarly derives falsity from its instantiated body. Concrete-symbol
\texttt{ho\_apply} introduction and elimination preserve the meaning of
predicate application.

Equality has reflexivity, symmetry, transitivity, and replacement rules.
Rewriting is an explicit derivation step, \texttt{Derives.equalSubst}:
\[
 \frac{\KB;\Gamma\vdash_N\operatorname{equal}(a,b)
       \qquad\KB;\Gamma\vdash_N F[a/z]}
      {\KB;\Gamma\vdash_N F[b/z]}.
\]
Both instances of the formula template $F$ use checked capture-avoiding
substitution. The context $\Gamma$ stays fixed. Replacement therefore reaches
inside operation arguments, dictionary expressions, and scoped goals while
respecting their binders. To rewrite using
$a\mathrel{\Leftrightarrow}b\mathbin{:-}\bar C$, first prove the instantiated
body $\bar C$ and derive the equality by clause application, then apply this
replacement rule. An oracle-certified equation uses the same two steps in the
extended KB. A resolver may automate these steps; the derivation records their
justification. The soundness proof for replacement uses equality of values and
the substitution lemma.

\texttt{Model.sound} includes these rules. Equations are symmetric; orienting
them for computation and proving rewriter termination are operational questions.
The checked example in \texttt{Tests/Foundation/EquationRules.lean} validates
$\operatorname{natEqual}(x,x)\Leftrightarrow\mathsf{true}$ under
$\operatorname{isNat}(x)$, then replaces a Boolean dictionary-call argument.

\section{Compositional models and rigid parameters}
A model $M$ has a nonempty value carrier $V$, interpretations of literals and
symbol applications, a truth predicate on values, and value-level logical
operations. Each logical operation satisfies its usual truth condition. For example,
\[
 M,\rho\models A\Rightarrow B
 \quad\Longleftrightarrow\quad
 (M,\rho\models A\Rightarrow M,\rho\models B),
\]
where $\rho:\mathrm{Variable}\to V$ is an environment. Quantifiers range over
$V$. Formula constructors also denote values, so formulas used as data retain
compositional semantics. The model includes the truth law for concrete-symbol
\texttt{ho\_apply}.
Its equality operation is true exactly when the operand values are equal;
having the same truth status alone does not imply equality of values.

Dictionary operation calls have the form
\texttt{dictionaryCall}$(d,\pi,f,\bar t)$, where $d$ and $\bar t$ are terms,
$\pi$ is a path through required-model subdictionaries, and $f$ is a field name.
Selectors use the symbol namespace. Dictionaries form finite trees of callable
leaves and subdictionaries; interfaces recursively require child interfaces and
their semantic laws. The checked interpretation is supplied by
\texttt{Extension.withDictionaryCalls} in \texttt{OperationDictionaries.lean}.
If lookup at $\pi.f$ accepts an operation entry with symbol $s$ and captures $\bar c$,
then
\[
 \lbrack\!\lbrack\operatorname{dictionaryCall}(d,\pi,f,\bar t)\rbrack\!\rbrack_{M,\rho}
 =\operatorname{application}_M
   (s,\bar c\mathbin{+\!+}\lbrack\!\lbrack\bar t\rbrack\!\rbrack_{M,\rho}).
\]
This is equality of returned values, allowing calls inside data expressions.
Invalid accesses denote an explicitly chosen value and supply no NAF evidence.
Substitution and opening traverse $d$ and every argument; the general soundness
proof covers these terms without adding inference rules.

\texttt{OrderDictionaries.lean} gives concrete \texttt{natOrd} and
\texttt{boolOrd} dictionaries in a shared untyped carrier, with $\mathsf{false}<\mathsf{true}$.
Each contains a WeakOrd child and a required Eq grandchild. The same term calls
\texttt{weakOrd.compare} through either dictionary. Interface proofs establish
exact domain membership, lawful equality, and a reflexive, transitive,
antisymmetric, total non-strict order. Comparison returns an \texttt{Ordering}
value; all laws are guarded by the instance's domain predicate.

A rigid parameter is a proof-level name, not a new element of $V$.
Its denotation is $\rho(a)$. The forall rule must work for every value assigned
to the fresh name:
\[
 \forall v\in V.\;M,\rho[a\mapsto v]\models F[a/x]
 \quad\Longrightarrow\quad M,\rho\models\forall x.F(x).
\]
Freshness ensures that changing $\rho(a)$ preserves the truth of $\Gamma$.
Syntactically distinct rigid names may denote the same value. In a pure Herbrand
interpretation, an environment assigns them ground terms; they are not added
to the universe as new constructors.

We write $M\models\KB$ when every clause is true under every environment,
and $M,\rho\models\Gamma$ when every local assumption holds under $\rho$.
The formal definition and the quantifier soundness argument are in
\texttt{Semantics.lean}.

\section{Inductive domains and validity of induction}
Let $\Sigma$ be a finite family of inductive domains defined by constructor
rules with finite field lists. The least generated membership relation is
$\operatorname{Member}_\Sigma(D,t)$ on finite constructor trees $t$.
Mutual recursion, empty domains, and overlap are permitted. Finitely many
definitions need not generate finitely many values.

Exact domain realization requires
\[
 M\models\operatorname{is}_D(v)
 \quad\Longleftrightarrow\quad
 \exists t.\;\operatorname{Member}_\Sigma(D,t)\land\lbrack\!\lbrack t\rbrack\!\rbrack_M=v.
\]
Here tree denotation is compositional. Constructor-closure clauses alone are
weaker: they would allow unsupported domain members. The implementation is
\texttt{InductiveDomains.Realizes}.

Define $\Dall(x)=\bigvee_{D\in\Sigma}\operatorname{is}_D(x)$. Then
\[
 \forall x:\Dall.\,P(x):=\forall x.\,\Dall(x)\Rightarrow P(x).
\]
This quantifies over denotations of generated members. Equivalence to raw forall
requires a proof that those denotations cover the entire semantic carrier.

For Nat, an induction principle can be an ordinary KB clause:
\[
\begin{split}
 \operatorname{apply}(P,x)\mathbin{:-}{}&\operatorname{Nat}(x),\\
 &\operatorname{apply}(P,\operatorname{zero}),\\
 &\forall n.\,\operatorname{Nat}(n)\Rightarrow
   (\operatorname{apply}(P,n)\Rightarrow
    \operatorname{apply}(P,\operatorname{succ}(n))).
\end{split}
\]
Apply denotes \texttt{ho\_apply}. To validate this clause, fix an arbitrary
environment and predicate value $P$. The base and step establish the invariant
on every generated natural by mathematical induction. Exact Nat membership
then establishes the head at any value satisfying the guard.

This validity proof is separate from using the clause in a derivation. In our
concrete model it is \texttt{InductionModel.inductionRule\_valid}. Generic
soundness uses the hypothesis $M\models\KB$, including validity of any induction
clauses present; it does not assert that arbitrary rules have a model.

\section{NAF evidence and adequacy}
NAF is its own syntax constructor. Its introduction rule is
\[
 \frac{\operatorname{FV}(G)=\varnothing\qquad N(\KB,\Gamma,G)}
      {\KB;\Gamma\vdash_N\naf G}.
\]
It is not justified by the mere absence of a derivation. Evidence is indexed
by the exact KB and actual local assumptions.

The general soundness theorem requires the failure provider to be adequate:
for every closed $G$ and every evidence $N(\KB,\Gamma,G)$,
\[
 \forall\rho.\;M,\rho\models\Gamma\Rightarrow M,\rho\models\naf G.
\]
This is \texttt{Model.FailureAdequate}. The general model requires coherence of
NAF truth with falsity of its operand. For the certified provider, we additionally
choose complement NAF:
\[
 M,\rho\models\naf G\quad\Longleftrightarrow\quad M,\rho\not\models G.
\]

The provider built from \texttt{CertifiedRefutation.Refutes} is proved adequate
under that interpretation. Its certificates establish semantic falsity in the
same original context. Atomic certificates require justified completion and
exhaustive coverage; implication refutation requires established antecedent
truth and consequent falsity. Guarded universal refutation exhibits an actual
generated member with a false property. Exhaustive counterexample-failure
certificates can also establish truth of a guarded universal, with separate
soundness lemmas for those semantic certificates.

These results do not prove that arbitrary unsuccessful SLD search supplies
adequate evidence. Unsupported, suspended, or truncated searches cannot be
used as refutation certificates.

\section{Soundness theorem and proof}
\begin{theorem}[Soundness of the logic]\label{thm:soundness}
Let $M$ be a compositional model, $M\models\KB$, and let $N$ be adequate for
$M$ and $\KB$. Then
\[
 \KB;\Gamma\vdash_N G
 \quad\Longrightarrow\quad
 \forall\rho.\;M,\rho\models\Gamma\Rightarrow M,\rho\models G.
\]
\end{theorem}

\begin{proof}[Proof outline]
Induct on the derivation in \texttt{Derives}.
The propositional rules follow from the truth laws of their connectives.
For clause application, compositional evaluation of a capture-avoiding
substitution agrees with evaluation in the composed environment. Clause validity
under every environment therefore transfers truth from the instantiated body
to its instantiated head.

For universal introduction, choose any $v\in V$ and interpret the fresh
parameter by $v$. Freshness preserves the assumptions, and the opening lemma
identifies the opened body with the quantified body at $v$. Quantifier
elimination and the existential rules similarly use substitution and freshness.
The application rules follow from the model's concrete-symbol application law.
Equality rules use equality of values and the substitution lemma, including
replacement inside scoped formulas and operation arguments.
NAF introduction follows from provider adequacy, which is an independent
obligation rather than a consequence of failed proof discovery.
\end{proof}

The general Lean theorem is \texttt{Semantics.Model.sound}. Exact inductive
realization, KB validity, and complement NAF are packaged in
\texttt{InductiveSoundness.Interpretation}. With the proved-adequate certified
provider, the logic theorem is:
\begin{quote}
\href{AnthillModKbSpec/Resolution/InductiveSoundness.lean}{\texttt{InductiveSoundness.declarative\_sound}}
\end{quote}
The theorem applies to the declarative logic independently of matching, clause
selection, backtracking, or execution strategy. Resolution proof translation
has its own soundness corollary; it is not a prerequisite for the logic theorem.

\begin{corollary}[Consistency]
If such a model exists, the positive calculus cannot derive falsity from the
empty local context. For a chosen model and its certified NAF provider, the
same holds with certified NAF enabled.
\end{corollary}
This follows because $M\not\models\bot$. The certified consistency theorem is:
\begin{quote}
\texttt{InductiveSoundness.declarative\_consistent}
\end{quote}

\subsection*{Extension by a sound equation oracle}
Let $O(\bar C,l,r)$ certify a conditional equation. Its soundness contract is
\[
 O(\bar C,l,r)\Longrightarrow
 \forall\rho.\;M,\rho\models\bigwedge\bar C\Longrightarrow
 \lbrack\!\lbrack l\rbrack\!\rbrack_{M,\rho}
 =\lbrack\!\lbrack r\rbrack\!\rbrack_{M,\rho}.
\]
Define $\KB_O$ as $\KB$ together with every conditional equation certified by $O$.
If $M\models\KB$, the contract establishes $M\models\KB_O$. Applying
Theorem~\ref{thm:soundness} proves soundness over $\KB_O$, provided the NAF
provider is adequate for this extended KB. Oracle acceptance retains the
equation's condition: the ordinary clause rule still proves that body.
Completeness is unnecessary, and a missing answer supplies no negative evidence.
The general theorem in \texttt{EquationOracle.lean} is:
\begin{quote}
\texttt{Model.sound\_with\_equation\_oracle}
\end{quote}
The inductive package preserves the same model, exact domain realization, and
complement NAF; its logic theorem is:
\begin{quote}
\texttt{InductiveSoundness.declarative\_sound\_with\_equation\_oracle}
\end{quote}
The corresponding consistency result excludes falsity in the empty context.
The checked oracle examples include a closed Nat computation, an equation used
in a predicate body, and a nonempty KB containing induction.

\section{A nonempty KB with a model}
The concrete witness in \texttt{InductionModel.lean} uses a carrier of constructor
keys paired with propositions, allowing both data and formulas as values.
It is Herbrand-style rather than a pure ground-term carrier. Nat membership is
exact and the denotations of different natural numbers are distinct.

The model satisfies Nat membership and induction clauses, followed by recursive
definitions of reachable and even predicates. The predicate interpretations
coincide with their least generated relations. In particular, evenness at one
is false, so this is not a model making every atom true. Model existence and
consistency are established for these nonempty example KBs, not for arbitrary
predicate extensions.

\section{Checking the formalization}
The project uses Lean 4.35.0. Check the logic and model results with:
\begin{verbatim}
lake build
lake env lean Tests/Resolution/InductiveSoundness.lean
lake env lean Tests/Models/InductionModel.lean
\end{verbatim}
The study currently covers soundness of the declarative calculus with adequate
NAF, a proved-adequate certified provider, and concrete inductive model witnesses.
It does not establish completeness or correctness of the production Anthill
resolver. Explicit negation and open/closed predicate policies remain future work.
\end{document}
